\chapter{Implementierung und Ausführung der QBF-Übersetzer}
\label{ch:appendix-qbf-tools}

Dieses Kapitel dokumentiert die Programme, mit denen Anfragen der Form
\(X\perp_\sigma Y\mid Z\) in quantifizierte boolesche Formeln
übersetzt und mit QuAbS ausgewertet werden. Berücksichtigt werden die
Semantiken
\(\sigma\in\{\mathrm{st},\mathrm{co},\mathrm{pr}\}\).

Die mathematische Herleitung der Kodierungen wurde in
Kapitel~\ref{ch:semantic-independence} behandelt. Im Mittelpunkt dieses
Kapitels stehen der Aufbau der Implementierung, die erwarteten
Ein- und Ausgaben sowie die reproduzierbare Ausführung der Programme.

Die Übersetzer erzeugen unmittelbar Dateien im schaltkreisbasierten
Format \texttt{QCIR-G14}. Die erzeugten Dateien können direkt an den 
QBF-Solver QuAbS übergeben werden.

\section{Werkzeuge und Einzelübersetzung}
\label{sec:appendix-qbf-single-translation}

\paragraph{Dateien und Verzeichnisstruktur.}

Die vier betrachteten Dateien und ihre Aufgaben sind in
Tabelle~\ref{tab:appendix-qbf-files} zusammengefasst.

\begin{table}[htb]
    \centering
    \begin{tabularx}{\textwidth}{@{}lX@{}}
        \toprule
        Datei & Aufgabe\\
        \midrule
        \texttt{af\_to\_stind\_qcir.py}
            & Übersetzt eine Anfrage zur
              \emph{stable}-Unabhängigkeit in \texttt{QCIR-G14}.\\
        \texttt{af\_to\_coind\_qcir.py}
            & Übersetzt eine Anfrage zur
              \emph{complete}-Unabhängigkeit in \texttt{QCIR-G14}.\\
        \texttt{af\_to\_prind\_qcir.py}
            & Übersetzt eine Anfrage zur
              \emph{preferred}-Unabhängigkeit einschließlich der
              Maximalitätsprüfung in \texttt{QCIR-G14}.\\
        \texttt{run\_batch\_singleton\_z.sh}
            & Führt die drei Übersetzer auf mehreren AFs und
              reproduzierbar gewählten Singleton-Anfragen aus, startet
              QuAbS und erfasst die Ergebnisse in CSV-Dateien.\\
        \bottomrule
    \end{tabularx}
    \caption{Dateien der QBF-basierten Auswertung}
    \label{tab:appendix-qbf-files}
\end{table}
Das Batchskript enthält zusätzlich eine optionale Phase zur Zählung von
Extensionen. Diese Phase gehört nicht zum hier dokumentierten
Dateiumfang. Deshalb wird sie in allen nachfolgenden Beispielaufrufen
mit \path{RUN_EXTENSION_COUNT=0} deaktiviert.

Bei Verwendung der im Batchskript hinterlegten Standardpfade wird die
folgende Verzeichnisstruktur erwartet:
\begin{lstlisting}[language={}]
Code/
|-- translate/
|   |-- af_to_stind_qcir.py
|   |-- af_to_coind_qcir.py
|   |-- af_to_prind_qcir.py
|   `-- run_batch_singleton_z.sh
|-- Test_Instances/
|   `-- final_curation/
|       `-- *.af
`-- quabs/
    `-- build/
        `-- quabs
\end{lstlisting}

Wird das Batchskript aus dem Verzeichnis
\texttt{Code/translate} gestartet, entstehen standardmäßig die
folgenden Ausgaben:
\begin{lstlisting}[language={}]
out_singleton_z/
|-- stable/
|-- complete/
`-- preferred/

results_singleton_z_stable.csv
results_singleton_z_complete.csv
results_singleton_z_preferred.csv
\end{lstlisting}

Die Eingabe-, Werkzeug- und Ausgabepfade können über
Umgebungsvariablen angepasst werden. Dabei ist zu beachten, dass die
Standardpfade der Ergebnisdateien und des Verzeichnisses
\texttt{out\_singleton\_z} relativ zum aktuellen Arbeitsverzeichnis
ausgewertet werden.

\paragraph{Voraussetzungen.}

Für die einzelnen Python-Übersetzer werden lediglich Python ab
Version~3.10 und die Python-Standardbibliothek benötigt. Die
Batchausführung setzt zusätzlich eine Linux- oder WSL-Umgebung mit
Bash ab Version~4.3 sowie üblichen GNU-Werkzeugen voraus. Hierzu zählen
insbesondere \texttt{awk}, \texttt{grep}, \texttt{sort},
\texttt{sha256sum}, \texttt{timeout} und \texttt{wc}. Bei paralleler
Ausführung wird außerdem \texttt{flock} benötigt.

Sofern nicht ausschließlich QCIR-Dateien erzeugt werden sollen, muss
eine ausführbare QuAbS-Binärdatei vorhanden sein. Ihr Pfad kann durch
die Umgebungsvariable \texttt{QUABS} festgelegt werden. Ohne explizite
Angabe verwendet das Batchskript
\texttt{../quabs/build/quabs} relativ zum übergeordneten
\texttt{Code}-Verzeichnis.


\paragraph{Eingabeformat.}

Die Übersetzer erwarten einen AF im ICCMA-Format. Die erste Zeile gibt
mit \texttt{p af \(n\)} die Anzahl der Argumente an. Optionale
Kommentarzeilen können anschließend die Argumentnamen in der
Reihenfolge ihrer 1-basierten Indizes festlegen. Eine Zeile
\texttt{\(i\) \(j\)} repräsentiert einen Angriff vom Argument \(i\) auf
das Argument \(j\).

\begin{Example}[ICCMA-Darstellung eines Zweizyklus]
\label{ex:appendix-qbf-af-input}
Der Zweizyklus
\[
F_{\leftrightarrow}
=
\bigl(
    \{a_1,a_2\},
    \{(a_1,a_2),(a_2,a_1)\}
\bigr)
\]
wird beispielsweise folgendermaßen dargestellt:
\begin{lstlisting}[
    language={},
    caption={AF im ICCMA-Format},
    label={lst:appendix-qbf-af-example}
]
p af 2
# a1
# a2
1 2
2 1
\end{lstlisting}
\end{Example}

Werden keine Namen angegeben, verwenden die Programme automatisch
\texttt{a1},\ldots,\texttt{a\(n\)}. Werden Namenskommentare verwendet,
sollten genau \(n\) eindeutige Namen angegeben werden.

Die Mengen \(X\), \(Y\) und \(Z\) werden über die Optionen
\texttt{-x}, \texttt{-y} und \texttt{-z} als kommaseparierte Listen
übergeben. Sie müssen paarweise disjunkt sein und dürfen intern keine
Duplikate enthalten. Die Programme akzeptieren sowohl Argumentnamen
als auch 1-basierte Indizes:
\begin{description}
    \item[\texttt{--argument-mode auto}]
        Namen werden bevorzugt; nicht als Name erkannte Eingaben
        werden als Indizes interpretiert. Dies ist der Standardmodus.
    \item[\texttt{--argument-mode index}]
        Sämtliche Angaben werden ausdrücklich als 1-basierte Indizes
        interpretiert. Das Batchskript verwendet stets diesen Modus.
\end{description}

\paragraph{Aufbau der Übersetzer und der QCIR-Ausgabe.}

Alle drei Programme folgen demselben grundlegenden Ablauf:
\begin{enumerate}
    \item Einlesen des AF und Aufbau der Angreifermengen,
    \item Einlesen und Validieren von \(X\), \(Y\) und \(Z\),
    \item Vergabe numerischer IDs für die quantifizierten Variablen,
    \item Erzeugung der semantikspezifischen Gatter sowie der
          Gleichheits- und Rekombinationsbedingungen,
    \item Ausgabe der QCIR-Datei. Sie umfasst Header und
          Quantorenpräfix sowie Ausgabegatter und Schaltkreismatrix.
\end{enumerate}
Die wesentlichen semantikspezifischen Unterschiede sind in
Tabelle~\ref{tab:appendix-qbf-encodings} zusammengefasst.

\begin{table}[htb]
    \centering
    \begin{tabularx}{\textwidth}{@{}lclX@{}}
        \toprule
        Semantik & Variablen & Präfix & Repräsentation\\
        \midrule
        \(\mathrm{st}\)
            & \(3n\)
            & \(\forall I^1\,\forall I^2\,\exists I^3\)
            & Eine In-Variable je Argument und Labellingkopie.\\
        \(\mathrm{co}\)
            & \(9n\)
            & \(\forall\mathbf V^1\,
               \forall\mathbf V^2\,
               \exists\mathbf V^3\)
            & Drei Variablen für die Labels
              \(\mathsf{in}\), \(\mathsf{out}\) und
              \(\mathsf{und}\).\\
        \(\mathrm{pr}\)
            & \(18n\)
            & \(\forall\mathbf V^1\,
               \forall\mathbf V^2\,
               \exists\mathbf V^3\,
               \exists\widehat{\mathbf V}^{1}\,
               \exists\widehat{\mathbf V}^{2}\,
               \forall\widehat{\mathbf V}^{3}\)
            & Drei Haupt- und drei Vergleichskopien zur Prüfung
              maximaler In-Mengen.\\
        \bottomrule
    \end{tabularx}
    \caption{Struktur der semantikspezifischen QCIR-Kodierungen}
    \label{tab:appendix-qbf-encodings}
\end{table}

Die QCIR-Ausgabe besteht jeweils aus
\begin{enumerate}
    \item der Kopfzeile \texttt{\#QCIR-G14} und kommentierten Metadaten,
    \item dem semantikspezifischen Quantorenpräfix,
    \item einer Anweisung \texttt{output(\(g\))} für das oberste Gatter
          und
    \item fortlaufend nummerierten \texttt{and}- und
          \texttt{or}-Gattern.
\end{enumerate}
Negative Gatter- oder Variablen-IDs bezeichnen Negationen.

\begin{Example}[Stable-Kodierung des Zweizyklus]
\label{ex:appendix-qbf-stable-encoding}
Als durchgängiges Beispiel dient der Zweizyklus
\(F_{\leftrightarrow}\) aus
Beispiel~\ref{ex:appendix-qbf-af-input}. Seine ICCMA-Darstellung ist in
Listing~\ref{lst:appendix-qbf-af-example} angegeben. Betrachtet wird die Anfrage
\(
    X=\{a_1\},
    Y=\{a_2\},
    Z=\emptyset.
\)
Die beiden \emph{stable}-Labellings des AFs sind durch
\(
    \mathsf{in}(\lambda_1)=\{a_1\}
    \text{ und }
    \mathsf{in}(\lambda_2)=\{a_2\}
\)
gegeben.

Da \emph{stable}-Labellings keine mit \(\mathsf{und}\) bezeichneten
Argumente enthalten, werden in der Implementierung ausschließlich
In-Variablen verwendet. Für jede Labellingkopie
\(I^k=\{i_a^k\mid a\in A\}\) lautet die verwendete Kodierung
\[
    \operatorname{St}^{I}_F(I^k)
    :=
    \bigwedge_{a\in A}
    \left(
        i_a^k
        \leftrightarrow
        \bigwedge_{b\in a^-}\neg i_b^k
    \right).
\]
Eine Variable \(i_a^k\) ist somit genau dann wahr, wenn keiner der
Angreifer von \(a\) in der zugehörigen In-Menge enthalten ist.

Für den Zweizyklus gilt daher
\[
\begin{split}
    \operatorname{St}^{I}_{F_{\leftrightarrow}}(I^k)
    =
    (i_{a_1}^k\leftrightarrow\neg i_{a_2}^k)
    \land
    (i_{a_2}^k\leftrightarrow\neg i_{a_1}^k).
\end{split}
\]
Da Gleichheitsbedingungen über der leeren Menge wahr sind, vereinfacht
sich die QBF der Anfrage zu
\[
\begin{split}
    \Phi^{\mathrm{st}}_{F_{\leftrightarrow}}
        (\{a_1\},\{a_2\}\mid\emptyset)
    :=
    \forall I^1\,\forall I^2\,\exists I^3\,
    \Bigl(
        &\operatorname{St}^{I}_{F_{\leftrightarrow}}(I^1)
        \land
        \operatorname{St}^{I}_{F_{\leftrightarrow}}(I^2)
        \\
        {}\rightarrow{}&
        \operatorname{St}^{I}_{F_{\leftrightarrow}}(I^3)
        \land
        (i_{a_1}^1\leftrightarrow i_{a_1}^3)
        \land
        (i_{a_2}^2\leftrightarrow i_{a_2}^3)
    \Bigr).
\end{split}
\]
Der Existenzquantor über \(I^3\) kann in das Präfix gezogen werden, da
dessen Variablen nicht in der Prämisse vorkommen.
\begin{table}[H]
    \centering
    \begin{tabularx}{\textwidth}{@{}clX@{}}
        \toprule
        QCIR-IDs & Variablen & Funktion\\
        \midrule
        \(1,2\)
            & \(i_{a_1}^1,i_{a_2}^1\)
            & Erste universell quantifizierte Labellingkopie.\\
        \(3,4\)
            & \(i_{a_1}^2,i_{a_2}^2\)
            & Zweite universell quantifizierte Labellingkopie.\\
        \(5,6\)
            & \(i_{a_1}^3,i_{a_2}^3\)
            & Existentiell quantifizierte Rekombination.\\
        \(7,\ldots,46\)
            & Hilfsgatter
            & Semantik-, Gleichheits- und Implikationsgatter.\\
        \bottomrule
    \end{tabularx}
    \caption{Variablenabbildung des Stable-Beispiels}
    \label{tab:appendix-qbf-stable-variables}
\end{table}
Die Zuordnung der QCIR-IDs ist in
Tabelle~\ref{tab:appendix-qbf-stable-variables} zusammengefasst.

Äquivalenzen werden gemäß
\(
    u\leftrightarrow v
    \equiv
    (\neg u\lor v)\land(u\lor\neg v)
\)
in \texttt{and}- und \texttt{or}-Gatter zerlegt. Die abschließende
Implikation wird als \(\neg A\lor B\) ausgegeben.
Listing~\ref{lst:appendix-qbf-stable-qcir-example} zeigt die dadurch
entstehende vollständige QuAbS-Eingabe.

\begin{lstlisting}[
    language={},
    caption={QCIR-G14-Eingabe des Stable-Beispiels},
    label={lst:appendix-qbf-stable-qcir-example}
]
#QCIR-G14
# AF source: beispiel.af
# |A| = 2
# X = {a1}
# Y = {a2}
# Z = {}
# Variablen-Mapping:
#   I^1: 1..2, I^2: 3..4, I^3: 5..6

forall(1, 2)
forall(3, 4)
exists(5, 6)
output(46)

7 = and(-2)
8 = or(-1, 7)
9 = or(1, -7)
10 = and(8, 9)
11 = and(-1)
12 = or(-2, 11)
13 = or(2, -11)
14 = and(12, 13)
15 = and(10, 14)
16 = and(-4)
17 = or(-3, 16)
18 = or(3, -16)
19 = and(17, 18)
20 = and(-3)
21 = or(-4, 20)
22 = or(4, -20)
23 = and(21, 22)
24 = and(19, 23)
25 = and(-6)
26 = or(-5, 25)
27 = or(5, -25)
28 = and(26, 27)
29 = and(-5)
30 = or(-6, 29)
31 = or(6, -29)
32 = and(30, 31)
33 = and(28, 32)
34 = and()
35 = or(-1, 5)
36 = or(1, -5)
37 = and(35, 36)
38 = and(37)
39 = or(-4, 6)
40 = or(4, -6)
41 = and(39, 40)
42 = and(41)
43 = and()
44 = and(15, 24, 34)
45 = and(33, 38, 42, 43)
46 = or(-44, 45)
\end{lstlisting}
Die drei Gattergruppen kodieren jeweils eine Kopie des
\emph{stable}-Labellings: \texttt{7} bis \texttt{15},
\texttt{16} bis \texttt{24} und \texttt{25} bis \texttt{33}.
Die Gatter \texttt{35} bis \texttt{38} erzwingen die Übereinstimmung
auf \(X\), die Gatter \texttt{39} bis \texttt{42} entsprechend auf
\(Y\). Die Gatter \texttt{34} und
\texttt{43} repräsentieren mit \texttt{and()} die wahre leere
Konjunktion für \(Z=\emptyset\). Schließlich bilden \texttt{44} die
Prämisse, \texttt{45} die Konklusion und \texttt{46} die Implikation.

Die entscheidende Ergebniszeile des Solveraufrufs lautet:
\begin{lstlisting}[language={}]
$QUABS out/beispiel_stable.qcir
r UNSAT
\end{lstlisting}
Zur Begründung können die universellen Blöcke die beiden
\emph{stable}-Labellings mit
\(
    I^1=(1,0)
    \text{ und }
    I^2=(0,1)
\)
kodieren. Die Gleichheitsbedingungen würden dann für \(I^3\) sowohl
\(i_{a_1}^3=1\) als auch \(i_{a_2}^3=1\) verlangen. Wegen des
gegenseitigen Angriffs ist \(\{a_1,a_2\}\) jedoch keine stabile
In-Menge. Es existiert somit keine passende Belegung von \(I^3\).
Die QBF ist falsch und entsprechend gilt
\(
    \{a_1\}
    \not\perp_{\mathrm{st}}
    \{a_2\}
    \mid\emptyset.
\)
\end{Example}

\paragraph{Einzelübersetzung und Solveraufruf.}

Zunächst wird das Zielverzeichnis für die QCIR-Dateien
angelegt:
\begin{lstlisting}[language={}]
mkdir -p out
\end{lstlisting}
Listing~\ref{lst:appendix-qbf-translation} zeigt die Aufrufe, mit denen für
\(X=\{a_1\}\), \(Y=\{a_2\}\) und \(Z=\emptyset\) die drei Kodierungen
erzeugt werden:
\begin{lstlisting}[
    language={},
    caption={Erzeugung der drei QCIR-Kodierungen},
    label={lst:appendix-qbf-translation}
]
python3 af_to_stind_qcir.py beispiel.af \
    --argument-mode index -x 1 -y 2 -z "" \
    -o out/beispiel_stable.qcir

python3 af_to_coind_qcir.py beispiel.af \
    --argument-mode index -x 1 -y 2 -z "" \
    -o out/beispiel_complete.qcir

python3 af_to_prind_qcir.py beispiel.af \
    --argument-mode index -x 1 -y 2 -z "" \
    -o out/beispiel_preferred.qcir
\end{lstlisting}
Ohne \texttt{-o} schreiben die Programme die erzeugte QCIR-Darstellung
auf die Standardausgabe. Bei Verwendung von \texttt{-o} muss das
übergeordnete Zielverzeichnis bereits existieren.

Eine erzeugte Datei wird QuAbS unmittelbar als Positionsargument
übergeben:
\begin{lstlisting}[language={}]
$QUABS out/beispiel_preferred.qcir
\end{lstlisting}
QuAbS gibt für eine wahre QBF \texttt{r SAT} und für eine falsche QBF
\texttt{r UNSAT} aus. In der vorliegenden Kodierung bedeutet
\texttt{SAT}, dass die untersuchte bedingte Unabhängigkeit gilt;
\texttt{UNSAT} bezeichnet entsprechend eine Abhängigkeit.

\section{Batchausführung und Ergebnisdateien}
\label{sec:appendix-qbf-batch-execution}

\paragraph{Batchablauf.}

Das Batchskript führt die drei Semantikphasen in der festen Reihenfolge
\emph{stable}, \emph{complete} und \emph{preferred} aus. Innerhalb
jeder Phase werden die ausgewählten AFs verarbeitet, bevor die nächste
Semantik beginnt.

Für jeden berücksichtigten AF führt das Skript die folgenden Schritte
aus:
\begin{enumerate}
    \item reproduzierbare Auswahl ungeordneter Singleton-Paare
          \(X=\{x\}\) und \(Y=\{y\}\),
    \item Erzeugung von sechs unabhängig gezogenen
          Konditionierungsmengen \(Z_0,\ldots,Z_5\),
    \item Übersetzung jeder Anfrage durch das zugehörige
          Python-Programm,
    \item Ausführung von QuAbS mit einem Zeitlimit und
    \item Erfassung von Anfrage, Formelgröße, Laufzeit und Ergebnis in
          der semantikspezifischen CSV-Datei.
\end{enumerate}
Standardmäßig werden fünf Argumentpaare je AF verwendet. Dabei gilt
\[
|Z_0|=0,
\qquad
|Z_5|=
\begin{cases}
\lfloor n/5\rfloor,
    &\text{für }\mathrm{st}\text{ und }\mathrm{co},\\
\lfloor n/10\rfloor,
    &\text{für }\mathrm{pr}.
\end{cases}
\]
Die vier Zwischengrößen werden gleichmäßig zwischen diesen Werten
verteilt. Die \(Z\)-Mengen werden für jedes Level neu gezogen und sind
daher im Allgemeinen nicht ineinander enthalten. Sie bleiben stets
disjunkt zu \(\{x,y\}\).

\paragraph{Aufruf und Konfiguration.}
Das Skript wird zweckmäßig aus seinem eigenen Verzeichnis gestartet:
\begin{lstlisting}[language={}]
cd Code/translate
chmod +x run_batch_singleton_z.sh
RUN_EXTENSION_COUNT=0 ./run_batch_singleton_z.sh
\end{lstlisting}
Ohne weiteres Argument verwendet das Skript den voreingestellten
AF-Ordner. Dieser lautet \path{../Test_Instances/final_curation}.
Ein abweichender Ordner kann als erstes Positionsargument angegeben werden:
\begin{lstlisting}[language={}]
RUN_EXTENSION_COUNT=0 \
./run_batch_singleton_z.sh /pfad/zu/af-dateien
\end{lstlisting}
Listing~\ref{lst:appendix-qbf-batch-smoke-test} zeigt einen kurzen
reproduzierbaren Funktionstest mit einer einzelnen kleinen Instanz:
\begin{lstlisting}[
    language={},
    caption={Kleiner reproduzierbarer Batchtest},
    label={lst:appendix-qbf-batch-smoke-test}
]
MIN_SIZE=2 LIMIT=1 NUM_PAIRS=1 KEEP_QCIR=1 \
RUN_EXTENSION_COUNT=0 \
./run_batch_singleton_z.sh /pfad/zu/af-dateien
\end{lstlisting}
Die explizite Angabe \texttt{MIN\_SIZE=2} ist erforderlich, weil der
Standardwert \(100\) beträgt und kleine Testinstanzen andernfalls
übersprungen werden.

Die zentralen Konfigurationsvariablen und ihre Standardwerte sind in
Tabelle~\ref{tab:appendix-qbf-batch-options} zusammengefasst.

\begin{table}[htb]
    \centering
    \begin{tabularx}{\textwidth}{@{}lcX@{}}
        \toprule
        Variable & Standard & Bedeutung\\
        \midrule
        \texttt{MIN\_SIZE} & \(100\)
            & Kleinste berücksichtigte Argumentzahl.\\
        \texttt{MAX\_SIZE} & \(1000\)
            & Größte berücksichtigte Argumentzahl.\\
        \texttt{LIMIT} & \(0\)
            & Maximale Zahl von AF-Dateien; \(0\) bedeutet keine
              Begrenzung.\\
        \texttt{NUM\_PAIRS} & \(5\)
            & Zahl der Singleton-Paare je AF.\\
        \texttt{SEED} & \(42\)
            & Ausgangswert der reproduzierbaren Zufallsauswahl.\\
        \texttt{TIMEOUT\_S} & \(600\)
            & Zeitlimit eines QuAbS-Aufrufs in Sekunden.\\
        \texttt{PARALLEL\_JOBS} & \(1\)
            & Maximale Zahl gleichzeitig bearbeiteter AFs innerhalb
              einer Semantikphase.\\
        \texttt{KEEP\_QCIR} & \(0\)
            & Bei \(1\) bleiben die erzeugten QCIR-Dateien erhalten.\\
        \texttt{TRANSLATE\_ONLY} & \(0\)
            & Bei \(1\) wird QuAbS nicht gestartet.\\
        \texttt{RESUME} & \(1\)
            & Bereits in der jeweiligen CSV enthaltene Anfragen werden
              übersprungen.\\
        \texttt{RUN\_STABLE} & \(1\)
            & Aktiviert die \emph{stable}-Phase.\\
        \texttt{RUN\_COMPLETE} & \(1\)
            & Aktiviert die \emph{complete}-Phase.\\
        \texttt{RUN\_PREFERRED} & \(1\)
            & Aktiviert die \emph{preferred}-Phase.\\
        \texttt{RUN\_EXTENSION\_COUNT} & \(1\)
            & Aktiviert die hier nicht betrachtete Extensionszählung;
              für den dokumentierten Dateiumfang ist \(0\) zu setzen.\\
        \bottomrule
    \end{tabularx}
    \caption{Zentrale Parameter des Batchskripts}
    \label{tab:appendix-qbf-batch-options}
\end{table}

\paragraph{Ergebnisdateien und Statuswerte.}
Für jede Semantik wird eine eigene CSV-Datei erzeugt. Ihre Spalten
lassen sich in fünf Gruppen einteilen:
\begin{description}
    \item[Instanz:]
        \texttt{instance}, \texttt{n} und \texttt{R};
    \item[Anfrage:]
        \texttt{pair\_index}, \texttt{x}, \texttt{y},
        \texttt{attack}, \texttt{z\_level}, \texttt{z\_size} und
        \texttt{z\_args};
    \item[Formelgröße:]
        \path{qbf_vars}, \path{qbf_gates} und \path{qbf_bytes};
    \item[Laufzeit:]
        \path{translate_seconds} und \path{solve_seconds};
    \item[Ergebnis:]
        \texttt{status}, \texttt{solver\_exit\_code},
        \texttt{solver\_message} und \texttt{independent}.
\end{description}
Die wichtigsten Statuswerte und ihre Bedeutung sind in
Tabelle~\ref{tab:appendix-qbf-status} zusammengefasst.
\begin{table}[htb]
    \centering
    \begin{tabularx}{\textwidth}{@{}lcX@{}}
        \toprule
        Status & \texttt{independent} & Bedeutung\\
        \midrule
        \texttt{SAT} & \texttt{yes}
            & Die kodierte bedingte Unabhängigkeit gilt.\\
        \texttt{UNSAT} & \texttt{no}
            & Die kodierte bedingte Unabhängigkeit gilt nicht.\\
        \texttt{TIMEOUT} & \texttt{?}
            & QuAbS überschritt das vorgegebene Zeitlimit.\\
        \texttt{TRANSLATE\_ERR} & \texttt{?}
            & Die QCIR-Datei konnte nicht erzeugt werden.\\
        \texttt{SOLVER\_ERR} & \texttt{?}
            & QuAbS lieferte kein auswertbares Ergebnis.\\
        \texttt{TRANSLATED} & \texttt{?}
            & Die Datei wurde mit
              \texttt{TRANSLATE\_ONLY=1} erzeugt, aber nicht gelöst.\\
        \bottomrule
    \end{tabularx}
    \caption{Statuswerte der Batchausführung}
    \label{tab:appendix-qbf-status}
\end{table}

QuAbS verwendet den Exitcode \(10\) für \texttt{SAT} und \(20\) für
\texttt{UNSAT}. Der durch das Zeitlimit erzeugte Exitcode \(124\) wird
als \texttt{TIMEOUT} erfasst.

\paragraph{Reproduzierbarkeit und Wiederaufnahme.}
Die Zufallsauswahl wird aus \texttt{SEED} und dem Instanznamen
abgeleitet und ist bei unveränderten Eingaben reproduzierbar. Bei einer
Änderung von Seed, Instanzen oder Studiendesign sollten neue
Ergebnisdateien verwendet werden. Die Resume-Prüfung identifiziert
bestehende Anfragen lediglich anhand von Instanz, \(x\), \(y\) und
\(Z\)-Level. \texttt{RESUME=0} allein leert eine vorhandene CSV-Datei
nicht, sondern kann zusätzliche Zeilen an die bestehende Datei
anhängen.
